\section{\ref{sec:recognition}장. 인식}

인식 속도, 첫 단계들의 구조, 응답의 선형 판독, 시간의 연속성에 의한 학습은 사람과 원숭이에서 측정되었다. 사슬 전체를 두 연산으로 환원하는 것과 영상에 의한 학습은 이 책의 모델로 확인했다. 착시의 설명은 근거가 탄탄한 것부터 논쟁 중인 것까지 있다.

{\footnotesize
\setlength{\tabcolsep}{3pt}\renewcommand{\arraystretch}{1.15}
\begin{longtable}{>{\raggedright}p{0.5cm}>{\raggedright}p{4.0cm}>{\raggedright}p{5.6cm}>{\raggedright}p{2.3cm}>{\raggedright\arraybackslash}p{1.7cm}}
\toprule
No. & 명제 & 실험과 측정한 것 & 출처 & 상태 \\
\midrule\endfirsthead
\toprule
No. & 명제 & 실험과 측정한 것 & 출처 & 상태 \\
\midrule\endhead
1 & 이동이 있을 때 픽셀 단위 견본 비교는 동전보다 낫지 않다; 단계 쌍~\eqref{eq:scstage}는 도형을 어디서나 인식한다 & \texttt{models/\allowbreak recognition\_models.py}, $24$개 점의 “망막”, $4000$번 시행: 뒤집힌 점의 비율 $0$, $0.1$, $0.2$에서 견본은 $0.49$, $0.51$, $0.52$; $S$와 $C$ 쌍은 $1.00$, $0.86$, $0.70$ & 이 장의 유도 & 모델 \\
2 & 그림 속 동물은 $\sim150$\,ms 만에 인식된다. 열 개쯤의 단계를 단계당 $10$--$20$\,ms씩 한 번 통과 & 사람, $20$\,ms 동안 보여 준 사진에서 “동물이 있는가” 과제: 두 답에 대한 유발 전위가 $\sim150$\,ms에 갈라진다. 원숭이: 첫 응답 시간이 단계마다 늘어난다(V1이 가장 먼저, 이어서 V2와 V4). 단계 수와 “단계당 스파이크 0개 또는 1개”는 이 수치에서 나온 추론 & \cite{Thorpe1996,Schmolesky1998} & 측정됨 \\
3 & $S$ 단계는 부분들의 AND, $C$ 단계는 위치들에 대한 최댓값이다(명제~\ref{prop:conveyor}) & 고양이, 일차 시각 피질: 단순 세포는 특정 위치의 특정 방위 가장자리에, 복합 세포는 더 넓은 영역의 같은 방위에 응답한다. 복합 세포의 세포 내 기록: 막대 두 개에 대한 응답은 많은 세포에서 두 응답 중 큰 쪽에 가깝고, 평균적으로 max 연산에 가장 가깝다. 상호 억제에 의한 최댓값은 명제~\ref{prop:wta}, 전체 활동에 의한 나눗셈은 총설 & \cite{HubelWiesel1962,Lampl2004,Carandini2012} & 측정됨 \\
4 & 사슬 위로 갈수록 조합은 커지고 복잡해지며, 응답은 위치에 점점 덜 좌우된다 & 원숭이, 자극을 “결정적 특징”까지 단순화: V4와 뒤쪽 하측두 피질 세포는 작은 수용장에서 형태, 색, 질감의 복잡한 조합에 응답한다; 앞쪽 하측두 피질에서는 수용장이 더 크다. 모델에서 $S$와 $C$의 교대가 이 그림을 재현한다 & \cite{KobatakeTanaka1994,Riesenhuber1999} & 측정됨 \\
5 & 합성곱 신경망은 이 교대를 되풀이하며 윗단계 응답을 가장 잘 예측한다; 물체는 뉴런 무리의 발화율에서 선형으로 읽힌다 & 인식을 학습했지만 뇌에 맞추지 않은 네트워크: 윗층은 원숭이 하측두 피질 뉴런의 응답을, 중간층은 V4의 응답을 예측한다. 무작위로 고른 세포 $\sim100$개의 활동에 대한 선형 분류기가 $12.5$\,ms부터의 창에서 위치와 크기가 달라도 물체와 범주를 알아본다 & \cite{Fukushima1980,Yamins2014,Hung2005} & 측정됨 \\
6 & 흔적을 가진 헵 규칙~\eqref{eq:traceinvariance}은 흔적이 $\sim20$\,ms보다 짧지 않으면 이름표 없는 영상으로 불변성을 가르친다 & 느린 특징과 흔적 규칙의 착상은 이론. \texttt{models/\allowbreak recognition\_models.py}, $C$ 뉴런 두 개, 입력 $16$개, $10$번 실행, 물체 사이 휴지 $0.3$\,s. $\tau_{\text{tr}}=5$\,ms에서 도형과의 일치는 평균 $0.52$, $10$\,ms에서 $0.56$; $20$, $50$, $200$\,ms에서 모든 실행이 $1.00$($30$\,ms에서는 열 번 중 한 번 틀림) & \cite{Foldiak1991,WiskottSejnowski2002} & 모델 \\
7 & 뇌에서 불변성은 시간의 연속성으로 학습된다 & 원숭이: 시야의 특정 위치로 눈이 도약하는 동안 물체를 몰래 바꿔치기했다; 하측두 피질 뉴런의 위치 불변성이 바꿔치기에 따라 변했고, 한 시간 만에 이미 유의했다. 이것이 시각 학습의 주된 규칙이라는 것은 가설 & \cite{LiDiCarlo2008} & 측정됨 \\
8 & 움직임: 방향 검출기, 이어서 넓은 영역에 대한 같은 AND/OR 쌍 & 토끼 망막의 방향 검출기; 원숭이 MT 영역에서 기록한 뉴런 $110$개 모두 방향 선택적이며, 움직임에 대한 응답이 V1보다 강하다 & \cite{Barlow1965,Albright1984} & 측정됨 \\
9 & 윗단계는 아래로 예측을 보내고, 위로는 오차가 올라간다 & 모델은 일차 시각 피질 수용장의 비고전적 효과를 설명한다; 개별 관찰은 부합하지만(총설) 파이프라인의 일반 구조로서는 증명되지 않음 & \cite{RaoBallard1999,Keller2018} & 가설 \\
10 & 어려운 그림은 되먹임과 여러 바퀴를 필요로 한다 & 원숭이: “어려운” 그림에서는 하측두 피질의 판정이 $\sim30$\,ms 늦게 형성된다; 이 늦은 응답은 순방향 네트워크가 잘 예측하지 못하고, 되먹임 네트워크나 더 깊은 네트워크가 더 잘 예측한다 & \cite{Kar2019} & 측정됨 \\
11 & 눈에 띄지 않는 픽셀 위조가 네트워크를 속인다; 사람의 시각은 훨씬 강하지만 무적은 아니다 & 네트워크: 특별히 고른 작은 섭동이 답을 바꾼다; “판다 $\to$ 긴팔원숭이” 예는 두 번째 논문에서. 사람에서는 짧게 보여 주면 네트워크 사이에 잘 전이되는 위조가 선택을 옮긴다 & \cite{Szegedy2014,Goodfellow2015,Elsayed2018} & 측정됨 \\
12 & 기대의 착시: 믿기 어려운 입력은 기대에 밀린다. 흐린 것은 “더 느리게” 움직이고, 드레스는 사람마다 다르게 보인다(\ref{sec:illusions}절) & 대비가 낮은 격자는 더 느리게 움직이는 것처럼 보인다. $1401$명 설문: 드레스는 세 가지 안정된 방식으로 보이고, 조명에 대한 힌트가 색을 바꾼다; 관찰자 $\sim13\,000$명에서 조명에 대한 가정이 결정한다. 느린 움직임을 기대하는 베이즈 모델이 여러 움직임 착시를 설명한다 & \cite{Thompson1982,LaferSousa2015,Wallisch2017,Weiss2002,Purves2011} & 가설 \\
13 & 이득 조절: 폭포, 기울기, 대비; 헤르만 격자는 이웃 억제가 아니다 & 토끼 망막의 방향 검출기는 오래 움직임을 보면 발화율이 떨어지고, 멈춘 뒤에는 배경 아래로 떨어진다. 기울기 착시는 이웃 활동에 의한 나눗셈 모델로 재현된다. 망막 출력 뉴런의 응답을 바꾸지 않는 격자 변형이 점을 없앤다. 이웃 억제 설명은 기각됨 & \cite{BarlowHill1963,Schwartz2009,SchillerCarvey2005} & 측정됨 \\
14 & 지연 보상: 움직이는 물체가 섬광보다 앞서 보인다 & 착시는 측정됨. 정신물리학은 움직임의 전방 연장과 지연 차이 둘 다와 모순된다; 섬광 뒤 $\sim80$\,ms의 사건에 따른 사후 설명이 제안되었다. 논쟁은 끝나지 않음 & \cite{Nijhawan1994,EaglemanSejnowski2000} & 가설 \\
15 & 정지한 “뱀”이 돈다: 지연 차이~\eqref{eq:latency}를 방향 검출기가 읽는다 & 착시는 색 없이도 일어난다; 이웃한 요소 쌍이 하나씩 착시를 만든다; 방향 선택적인 원숭이 피질 뉴런은 이 정지한 쌍에 방향성 응답을 낸다. 사람에서 회전은 미세 눈 도약과 눈 깜박임이 촉발한다 & \cite{Conway2005,OteroMillan2012} & 측정됨 \\
16 & 이중 그림: 상호 억제, 피로, 잡음; 전환은 주로 잡음이 일으킨다 & 경쟁하는 뉴런 무리 모델: 전환은 \emph{잡음}이 일으키며, 잡음이 없으면 전환도 없다; 자극 세기에 따른 전환 빈도 증가를 설명한다. 여기서 피로의 역할은 작다 & \cite{MorenoBote2007} & 가설 \\
17 & 착시의 일부는 기계가 배우는 과제의 결과이다 & 영상의 다음 프레임을 예측하도록 학습한 네트워크는 정지한 고리의 회전을 “보고” 대조 그림에서는 보지 않는다; 이미지의 잡음 제거와 선명화를 위한 네트워크는 사람처럼 색과 밝기 착시에 넘어간다 & \cite{Watanabe2018,GomezVilla2020} & 측정됨 \\
\bottomrule
\end{longtable}}
